\subsection{Statement of the main theorem}

Let X be a set. If \(m \leqslant n\) , we shall denote by \(\pi_{m,n}\) the map from \(X^{n}\) into \(X^{m}\) such that \(\pi_{m,n}(\boldsymbol{x}) = (x_{1}, \ldots, x_{m})\) for all \(\boldsymbol{x} = (x_{1}, \ldots, x_{n}) \in \mathrm{X}^{n}\) .

Let A be a unital Banach algebra over C. For every integer \(n \geqslant 1\) and every \(a \in A^{n}\) , we denote by \(\mathrm{Sp}^{n}(\boldsymbol{a})\) the simultaneous spectrum \(\mathrm{Sp}_{\mathrm{A}}^{\{1,\ldots,n\}}(\boldsymbol{a})\) (def. 2 of I, p.~42). It is a compact subset of \(C^{n}\) . For every integer m such that \(1 \leqslant m \leqslant n\) , one has \(\pi_{m,n}(\mathrm{Sp}^{n}(\boldsymbol{a})) = \mathrm{Sp}^{m}(\pi_{m,n}(\boldsymbol{a}))\) (I, p.~41, no. 6). The continuous linear map

\[
\pi_ {m, n} ^ {*} \colon \mathcal {O} (\mathrm{Sp} ^ {m} (\pi_ {m, n} (\boldsymbol {a})); \mathrm{A}) \longrightarrow \mathcal {O} (\mathrm{Sp} ^ {n} (\boldsymbol {a}); \mathrm{A})
\]

is a morphism of unital algebras.

Let A be a commutative unital Banach algebra over C. Let \(n \geqslant 1\) be an integer. One calls holomorphic functional calculus in n variables on A the data, for every \(a \in A^{n}\) , of a map

\[
\Theta_ {\boldsymbol {a}} \colon \mathcal {O} (\mathrm{Sp} ^ {n} (\boldsymbol {a}); \mathrm{A}) \longrightarrow \mathrm{A}
\]

satisfying the conditions:

(CF1) For every \(a \in A^{n}\) , the map \(\Theta_{a}\) is a continuous morphism of unital algebras.

(CF2) If \(\boldsymbol{a}=(a_{1},\ldots,a_{n})\) , and if \(z_{1},\ldots,z_{n}\) denote the germs in a neighborhood of \(\mathrm{Sp}^{n}(\boldsymbol{a})\) of the coordinate functions on \(C^{n}\) , one has

\[
\Theta_ {\boldsymbol {a}} (z _ {1}) = a _ {1}, \dots , \Theta_ {\boldsymbol {a}} (z _ {n}) = a _ {n}.
\]

Remark. --- If the radical of the algebra A is zero, one may omit the continuity condition in (CF1) (cf.~prop. 9 of I, p.~40).

One calls holomorphic functional calculus on A the data, for every integer \(n \geqslant 1\) , of a holomorphic functional calculus in n variables on A, satisfying:

(CF3) Whatever the integers m and n such that \(1 \leqslant m \leqslant n\) , and whatever \(a \in A^{n}\) and \(f \in \mathcal{O}(\mathrm{Sp}^{m}(\pi_{m,n}(a); \mathrm{A})\) , one has

\[
\Theta_ {\boldsymbol {a}} (\pi_ {m, n} ^ {*} (f)) = \Theta_ {\pi_ {m, n} (\boldsymbol {a})} (f).
\]

The object of this paragraph is to prove the following theorem:

THEOREM 1. --- Let A be a commutative complex unital Banach algebra. There exists a unique holomorphic functional calculus on A.

The proof of this theorem will occupy Nos. 3 to 7.

